\documentclass[UTF8,oneside,12pt]{ctexbook}
\usepackage[a4paper,top=2.45cm,bottom=2.55cm,left=2.35cm,right=2.35cm]{geometry}
\usepackage{fontspec}
\setmainfont{Liberation Serif}
\setsansfont{DejaVu Sans}
\setmonofont{DejaVu Sans Mono}
\setCJKmainfont{Noto Serif CJK SC}
\setCJKsansfont{Noto Sans CJK SC}
\setCJKmonofont{Noto Sans CJK SC}
\usepackage{amsmath,amssymb,amsthm,mathtools}
\usepackage{booktabs,longtable,array,multirow}
\usepackage{xcolor}
\usepackage{graphicx}
\usepackage{tikz}
\usetikzlibrary{arrows.meta,positioning,calc,fit,shapes.geometric,shapes.misc,trees}
\usepackage{listings}
\usepackage{hyperref}
\hypersetup{colorlinks=true,linkcolor=blue!40!black,urlcolor=blue!50!black,citecolor=blue!40!black,pdftitle={图与网络算法及树结构复习讲义},pdfauthor={ChatGPT}}
\linespread{1.18}
\setlength{\parindent}{2em}
\setlength{\parskip}{0.25em}
\definecolor{notegray}{RGB}{245,247,250}
\definecolor{noteblue}{RGB}{41,98,159}
\definecolor{darkred}{RGB}{160,0,0}
\definecolor{darkgreen}{RGB}{0,110,70}
\lstdefinestyle{pystyle}{language=Python,basicstyle=\ttfamily\footnotesize,keywordstyle=\color{blue!60!black},commentstyle=\color{darkgreen},stringstyle=\color{darkred},numbers=left,numberstyle=\tiny\color{gray},numbersep=7pt,showstringspaces=false,breaklines=true,frame=single,rulecolor=\color{gray!55},tabsize=4}
\lstset{style=pystyle}
\newtheorem{definition}{定义}[chapter]
\newtheorem{theorem}{定理}[chapter]
\newtheorem{lemma}{引理}[chapter]
\newtheorem{example}{例}[chapter]
\newenvironment{summarybox}{\begin{quote}\small\noindent\textbf{本节要点}\quad}{\end{quote}}
\newcommand{\algo}[1]{\textsf{#1}}
\newcommand{\bigO}{\mathcal O}
\title{图与网络算法及树结构\\\large 教学书本式复习讲义（20-30页版）}
\author{按“导言-理论-案例-步骤/代码-结论”结构整理}
\date{\today}
\begin{document}
\maketitle
\frontmatter
\chapter*{使用说明}
这份讲义把图论与网络算法、组合优化方法、常见平衡树和编码树放在同一条主线上。每个主题都尽量按照教材写法组织：先说明“解决什么问题”，再给出关键思想、理论依据、简单例题、可执行步骤或 Python 参考实现，最后说明适用条件与常见错误。篇幅有限，所以对证明采用“考试可复述”的深度；对最大流、最小费用最大流、A*、Johnson、红黑树等细节较多的方法，主要给出核心定理和实现框架。\\
编译建议：使用 XeLaTeX 编译本文件。本文使用 ctexbook、TikZ、booktabs、longtable 与 listings。中文排版采用 ctex，图示采用 TikZ，表格采用 booktabs/longtable，代码采用 listings。
\tableofcontents
\mainmatter

\chapter{图模型、遍历与有向无环图}
\section{图模型的基本语言}
\begin{definition}[图、路径与权]
图可记为 $G=(V,E)$，其中 $V$ 是顶点集，$E$ 是边集。若边有方向，则称有向图；若每条边 $(u,v)$ 还有数值 $w(u,v)$，则称带权图。路径是顶点序列 $v_0,v_1,\ldots,v_k$，且相邻顶点之间有边连接；路径长度可按边数计算，也可按权值和计算。
\end{definition}

图算法的第一步不是写代码，而是判断问题属于哪一类模型：
\begin{itemize}
  \item 只问“能不能到达、是否连通、有没有环”，通常从 DFS/BFS 出发；
  \item 边权全非负，求单源最短路，优先考虑 Dijkstra；
  \item 允许负边但不允许负环，考虑 Bellman-Ford 或 Johnson；
  \item 要求连接所有点且总代价最小，考虑最小生成树；
  \item 有容量、费用、供需、匹配关系时，考虑网络流或二分图匹配；
  \item 遍历所有城市、图着色、排课等常常是 NP-困难问题，需区分精确算法和启发式方法。
\end{itemize}

\begin{figure}[htbp]
\centering
\begin{tikzpicture}[node distance=9mm,>=Latex,box/.style={rectangle,draw,rounded corners,align=center,minimum width=25mm,minimum height=7mm,font=\small}]
\node[box] (g) {图模型};
\node[box,above right=of g] (tr) {DFS/BFS};
\node[box,right=of tr] (dag) {拓扑排序\\关键路径};
\node[box,right=of g] (sp) {最短路};
\node[box,right=of sp] (dij) {Dijkstra\\Bellman-Ford\\Floyd/Johnson/A*};
\node[box,below right=of g] (mst) {最小生成树};
\node[box,right=of mst] (kp) {Kruskal/Prim};
\node[box,below=of mst] (flow) {网络流与匹配};
\node[box,right=of flow] (ff) {最大流最小割\\最小费用流\\二分图匹配};
\node[box,below=of flow] (hard) {组合优化};
\node[box,right=of hard] (meta) {TSP/图着色\\GA/SA/ACO};
\node[box,below=of hard] (tree) {树结构};
\node[box,right=of tree] (trees) {B树/B+树\\红黑树/赫夫曼树};
\draw[->] (g)--(tr);\draw[->](tr)--(dag);\draw[->](g)--(sp);\draw[->](sp)--(dij);\draw[->](g)--(mst);\draw[->](mst)--(kp);\draw[->](g)--(flow);\draw[->](flow)--(ff);\draw[->](g)--(hard);\draw[->](hard)--(meta);\draw[->](g)--(tree);\draw[->](tree)--(trees);
\end{tikzpicture}
\caption{本讲义的知识主线}
\end{figure}

\section{DFS 与 BFS}
\subsection{导言与理论}
深度优先搜索 DFS 像“沿一条路走到底再回退”，适合发现连通分量、环、拓扑序、割点、桥等结构；广度优先搜索 BFS 像“逐层扩散”，适合无权最短路、层次遍历和二分图判定。两者在邻接表上都只访问每个顶点和每条边常数次，因此复杂度为
\[
\bigO(|V|+|E|).
\]

\begin{theorem}[BFS 的无权最短路性质]
在无权图中，从源点 $s$ 进行 BFS，第一次访问到顶点 $v$ 时得到的层数，等于从 $s$ 到 $v$ 的最短边数。
\end{theorem}
\begin{proof}
BFS 按层扩展：第 $k$ 层全部由第 $k-1$ 层顶点一步到达。因此若某点第一次在第 $k$ 层出现，必存在长度为 $k$ 的路径。若还有更短路径，则它应在更早层被发现，矛盾。
\end{proof}

\begin{figure}[htbp]
\centering
\begin{tikzpicture}[>=Latex,vertex/.style={circle,draw,minimum size=8mm}]
\node[vertex] (A) at (0,0) {A};
\node[vertex] (B) at (2,1) {B};
\node[vertex] (C) at (2,-1) {C};
\node[vertex] (D) at (4,1.2) {D};
\node[vertex] (E) at (4,0) {E};
\node[vertex] (F) at (4,-1.2) {F};
\draw (A)--(B)--(D);\draw (A)--(C)--(F);\draw (B)--(E);\draw (C)--(E);\draw (E)--(F);
\node[draw=none,below=8mm of A,align=left] {从 A 出发 BFS 层次：\\第 0 层 A；第 1 层 B,C；第 2 层 D,E,F。};
\end{tikzpicture}
\caption{BFS 逐层扩展示意}
\end{figure}

\subsection{Python 参考实现}
\begin{lstlisting}
from collections import deque

def bfs(graph, s):
    dist = {s: 0}
    parent = {s: None}
    q = deque([s])
    while q:
        u = q.popleft()
        for v in graph.get(u, []):
            if v not in dist:
                dist[v] = dist[u] + 1
                parent[v] = u
                q.append(v)
    return dist, parent

def dfs(graph, s):
    seen, order = set(), []
    def visit(u):
        seen.add(u)
        order.append(u)
        for v in graph.get(u, []):
            if v not in seen:
                visit(v)
    visit(s)
    return order
\end{lstlisting}

\section{拓扑排序与关键路径法}
\subsection{拓扑排序}
拓扑排序只适用于有向无环图 DAG。若边 $u\to v$ 表示“任务 $u$ 必须先于任务 $v$”，则拓扑序就是一种合法执行顺序。

Kahn 算法的思想是：反复取入度为 $0$ 的点输出，并删除它的出边。若最后输出点数少于 $|V|$，说明图中有环。
\[
\text{复杂度}=\bigO(|V|+|E|).
\]

\subsection{关键路径法 CPM}
关键路径法把项目看成 DAG，边或点表示活动，权值表示活动持续时间。求项目总工期，就是求从开始点到结束点的最长路径；在 DAG 中，最长路可按拓扑序动态规划。

设活动网络中 $ve[v]$ 表示事件 $v$ 的最早发生时间，则
\[
ve[v]=\max_{(u,v)\in E}\{ve[u]+w(u,v)\}.
\]
反向计算最迟发生时间 $vl[v]$ 后，活动 $(u,v)$ 的时间余量为
\[
slack(u,v)=vl[v]-ve[u]-w(u,v).
\]
余量为 $0$ 的活动就是关键活动，所有关键活动形成的路径称为关键路径。

\begin{example}[简单项目网络]
若活动 $A\to B$ 需 3 天，$A\to C$ 需 2 天，$B\to D$ 需 4 天，$C\to D$ 需 6 天，则两条路径长度分别为 $7$ 与 $8$，故关键路径为 $A\to C\to D$，项目最短工期为 $8$ 天。
\end{example}

\begin{summarybox}
DFS 偏结构分析，BFS 偏层次和无权最短路；拓扑排序必须要求 DAG；关键路径法本质上是在 DAG 上求最长路。考试中经常要求：写遍历序、判断有无环、给拓扑序、算最早/最迟时间和关键活动。
\end{summarybox}

\chapter{最短路算法族}
\section{Dijkstra 算法}
\subsection{导言}
Dijkstra 算法解决非负权图上的单源最短路径问题。给定源点 $s$，求它到所有顶点的最短距离。它是典型的贪心算法：每一步把当前暂定距离最小的顶点“定型”。

\subsection{核心理论}
\begin{theorem}[Dijkstra 定型正确性]
若所有边权非负，则每次从未定型顶点中选出 $d[u]$ 最小的顶点 $u$ 时，有 $d[u]=\delta(s,u)$。
\end{theorem}
\begin{proof}
反设存在一条更短的 $s\leadsto u$ 路径。在这条路径上取第一个未定型顶点 $x$，其前驱 $y$ 已定型。由于 $y$ 定型时会松弛边 $(y,x)$，可得 $d[x]$ 不大于真实最短距离。又因边权非负，从 $x$ 到 $u$ 不会减少长度，所以 $d[x]<d[u]$，与 $u$ 是未定型点中距离最小者矛盾。
\end{proof}

\begin{figure}[htbp]
\centering
\begin{tikzpicture}[>=Latex,vertex/.style={circle,draw,minimum size=8mm}]
\node[vertex] (s) at (0,0) {$s$};
\node[vertex] (a) at (2,1.3) {$a$};
\node[vertex] (b) at (2,-1.3) {$b$};
\node[vertex] (c) at (4,1.3) {$c$};
\node[vertex] (d) at (4,-1.3) {$d$};
\node[vertex] (t) at (6,0) {$t$};
\draw[->] (s)--node[above left,draw=none]{2}(a);
\draw[->] (s)--node[below left,draw=none]{5}(b);
\draw[->] (a)--node[above,draw=none]{2}(c);
\draw[->] (a)--node[left,draw=none]{1}(b);
\draw[->] (a)--node[right,draw=none]{4}(d);
\draw[->] (b)--node[below,draw=none]{1}(d);
\draw[->] (c)--node[above right,draw=none]{3}(t);
\draw[->] (d)--node[below right,draw=none]{2}(t);
\end{tikzpicture}
\caption{Dijkstra 示例图}
\end{figure}

从 $s$ 出发，距离变化为：初始 $d[s]=0$；选 $s$ 后 $d[a]=2,d[b]=5$；选 $a$ 后 $d[b]=3,d[c]=4,d[d]=6$；选 $b$ 后 $d[d]=4$；选 $c$ 后 $d[t]=7$；选 $d$ 后 $d[t]=6$。最终最短路 $s\to a\to b\to d\to t$，长度为 6。

\begin{lstlisting}
import heapq
from math import inf

def dijkstra(graph, source):
    dist = {v: inf for v in graph}
    prev = {v: None for v in graph}
    dist[source] = 0
    heap = [(0, source)]
    while heap:
        du, u = heapq.heappop(heap)
        if du != dist[u]:
            continue
        for v, w in graph[u]:
            if dist[v] > du + w:
                dist[v] = du + w
                prev[v] = u
                heapq.heappush(heap, (dist[v], v))
    return dist, prev
\end{lstlisting}

邻接表加二叉堆的常见复杂度为 $\bigO((V+E)\log V)$。注意：只要出现负权边，就不能直接使用 Dijkstra。

\section{Bellman-Ford 算法}
Bellman-Ford 的核心是反复松弛所有边。经过第 $k$ 轮后，所有“边数不超过 $k$ 的最短路”都已被正确考虑。由于没有负环时，简单最短路最多包含 $|V|-1$ 条边，所以做 $|V|-1$ 轮即可。

\begin{theorem}
若从源点可达的部分不存在负权回路，则 Bellman-Ford 在 $|V|-1$ 轮松弛后得到所有最短距离；若第 $|V|$ 轮仍能松弛某条边，则存在源点可达负环。
\end{theorem}

\begin{lstlisting}
def bellman_ford(vertices, edges, source):
    INF = 10**18
    dist = {v: INF for v in vertices}
    dist[source] = 0
    for _ in range(len(vertices) - 1):
        changed = False
        for u, v, w in edges:
            if dist[u] != INF and dist[v] > dist[u] + w:
                dist[v] = dist[u] + w
                changed = True
        if not changed:
            break
    for u, v, w in edges:
        if dist[u] != INF and dist[v] > dist[u] + w:
            raise ValueError("negative cycle reachable")
    return dist
\end{lstlisting}
复杂度为 $\bigO(VE)$，优点是能处理负边并检测负环。

\section{Floyd-Warshall 算法}
Floyd-Warshall 是全源最短路动态规划。设 $d_k[i][j]$ 表示从 $i$ 到 $j$ 的路径中，只允许使用编号不超过 $k$ 的顶点作为中间点时的最短距离，则转移为
\[
d_k[i][j]=\min\{d_{k-1}[i][j],\ d_{k-1}[i][k]+d_{k-1}[k][j]\}.
\]
实现时可把三维数组压成二维数组，外层枚举中间点 $k$：
\begin{lstlisting}
def floyd_warshall(dist):
    # dist[i][j] is initial edge weight, INF if no edge
    n = len(dist)
    for k in range(n):
        for i in range(n):
            for j in range(n):
                if dist[i][j] > dist[i][k] + dist[k][j]:
                    dist[i][j] = dist[i][k] + dist[k][j]
    return dist
\end{lstlisting}
复杂度为 $\bigO(V^3)$，适合点数不大、图较稠密、需要任意两点最短路的场景。

\section{Johnson 算法}
Johnson 用于稀疏图全源最短路，允许负边但不允许负环。它的关键是“重新赋权”：先加超级源点 $q$，用 Bellman-Ford 求势函数 $h(v)=\delta(q,v)$，再把每条边权改成
\[
w'(u,v)=w(u,v)+h(u)-h(v).
\]
由三角不等式可得 $w'(u,v)\ge 0$，所以之后可以从每个源点运行 Dijkstra。重赋权不会改变路径之间的相对大小，因为任一路径 $p:s\leadsto t$ 有
\[
w'(p)=w(p)+h(s)-h(t),
\]
对固定的 $s,t$，后半部分是常数。

\section{A* 搜索}
A* 可看成带启发函数的 Dijkstra。每个候选点用
\[
f(n)=g(n)+h(n)
\]
排序，其中 $g(n)$ 是起点到当前点的已知代价，$h(n)$ 是当前点到目标点的估计代价。若 $h$ 从不高估真实距离，称为 admissible；若还满足三角不等式形式，称为 consistent。启发函数越接近真实距离，搜索越少；若 $h\equiv0$，A* 退化为 Dijkstra。

\begin{summarybox}
最短路选择口诀：无权图用 BFS；非负权单源用 Dijkstra；有负边单源用 Bellman-Ford；点数较少的全源用 Floyd-Warshall；稀疏图全源且可有负边用 Johnson；有明确目标并有好启发函数时用 A*。
\end{summarybox}

\chapter{最小生成树、网络流与匹配}
\section{最小生成树：Kruskal 与 Prim}
\subsection{问题定义}
在连通无向带权图中，生成树连接所有顶点且无环。最小生成树 MST 是权值和最小的生成树。

\begin{lemma}[割性质]
对任意割 $(S,V\setminus S)$，跨越该割的最小权边必属于某棵最小生成树。
\end{lemma}
\begin{proof}
设 $e$ 是跨割最小边，任取一棵 MST $T$。若 $e\in T$ 则成立。若 $e\notin T$，把 $e$ 加入 $T$ 会形成一个环，环上必有另一条跨越该割的边 $f$。由 $w(e)\le w(f)$，用 $e$ 替换 $f$ 后仍是生成树且总权不增，因此存在包含 $e$ 的 MST。
\end{proof}

\subsection{Kruskal}
Kruskal 把所有边按权升序排列，从小到大尝试加入，只要不形成环就接受。判环用并查集。

\begin{lstlisting}
class DSU:
    def __init__(self, n):
        self.parent = list(range(n))
        self.rank = [0] * n
    def find(self, x):
        while self.parent[x] != x:
            self.parent[x] = self.parent[self.parent[x]]
            x = self.parent[x]
        return x
    def union(self, a, b):
        ra, rb = self.find(a), self.find(b)
        if ra == rb:
            return False
        if self.rank[ra] < self.rank[rb]:
            ra, rb = rb, ra
        self.parent[rb] = ra
        if self.rank[ra] == self.rank[rb]:
            self.rank[ra] += 1
        return True

def kruskal(n, edges):
    dsu = DSU(n)
    ans = []
    for w, u, v in sorted(edges):
        if dsu.union(u, v):
            ans.append((u, v, w))
    return ans
\end{lstlisting}
复杂度主要由排序决定，为 $\bigO(E\log E)$。Kruskal 更适合稀疏图。

\subsection{Prim}
Prim 从一个顶点开始，不断选择“当前树到外部顶点”的最小边。它与 Dijkstra 形式相似，但含义不同：Dijkstra 维护源点到顶点的路径长度，Prim 维护顶点接入当前生成树的最便宜边。

\begin{figure}[htbp]
\centering
\begin{tikzpicture}[vertex/.style={circle,draw,minimum size=8mm},>=Latex]
\node[vertex] (A) at (0,0) {A};
\node[vertex] (B) at (2,1.2) {B};
\node[vertex] (C) at (2,-1.2) {C};
\node[vertex] (D) at (4,1.2) {D};
\node[vertex] (E) at (4,-1.2) {E};
\draw[very thick] (A)--node[above left,draw=none]{1}(B);
\draw (A)--node[below left,draw=none]{4}(C);
\draw[very thick] (B)--node[above,draw=none]{2}(D);
\draw[very thick] (C)--node[below,draw=none]{3}(E);
\draw[very thick] (B)--node[left,draw=none]{2}(C);
\draw (D)--node[right,draw=none]{5}(E);
\draw (B)--node[above right,draw=none]{6}(E);
\end{tikzpicture}
\caption{粗线表示一棵可能的最小生成树}
\end{figure}

\section{最大流与最小割}
\subsection{模型}
网络流给定源点 $s$、汇点 $t$ 和容量 $c(u,v)$。流量 $f(u,v)$ 需满足容量限制 $0\le f(u,v)\le c(u,v)$ 与流量守恒。目标是最大化从 $s$ 到 $t$ 的总流量。

割 $(S,T)$ 满足 $s\in S,t\in T$，割容量为所有从 $S$ 指向 $T$ 的边容量和：
\[
c(S,T)=\sum_{u\in S,v\in T}c(u,v).
\]

\begin{theorem}[最大流最小割定理]
任一网络中，最大 $s-t$ 流的值等于最小 $s-t$ 割的容量。
\end{theorem}

增广路方法维护残量网络：若残量网络中还有从 $s$ 到 $t$ 的路，就沿路增加流量。Edmonds-Karp 规定每次用 BFS 找最短增广路，复杂度为 $\bigO(VE^2)$；Dinic 用分层图和阻塞流，实践中更快。

\begin{figure}[htbp]
\centering
\begin{tikzpicture}[>=Latex,vertex/.style={circle,draw,minimum size=8mm}]
\node[vertex] (s) at (0,0) {$s$};
\node[vertex] (a) at (2,1.1) {$a$};
\node[vertex] (b) at (2,-1.1) {$b$};
\node[vertex] (c) at (4,1.1) {$c$};
\node[vertex] (d) at (4,-1.1) {$d$};
\node[vertex] (t) at (6,0) {$t$};
\draw[->] (s)--node[above left,draw=none]{10}(a);
\draw[->] (s)--node[below left,draw=none]{8}(b);
\draw[->] (a)--node[above,draw=none]{5}(c);
\draw[->] (a)--node[left,draw=none]{4}(b);
\draw[->] (b)--node[below,draw=none]{10}(d);
\draw[->] (c)--node[right,draw=none]{7}(t);
\draw[->] (d)--node[below right,draw=none]{10}(t);
\draw[->] (c)--node[above right,draw=none]{3}(d);
\draw[dashed,red!70!black,thick] (3,-1.8)--(3,1.8);
\node[draw=none,red!70!black] at (3,1.95) {一个割};
\end{tikzpicture}
\caption{网络流与割的直观关系}
\end{figure}

\section{最小费用最大流}
最小费用最大流在每条边上同时有容量 $c(u,v)$ 与单位费用 $a(u,v)$。目标是在最大流量的前提下，使总费用
\[
\sum_{(u,v)\in E}a(u,v)f(u,v)
\]
最小。常见入门做法是“逐次最短增广路”SSAP：每次在残量网络中找一条从 $s$ 到 $t$ 的最小费用路径，再沿路增广。若有负费用边，常结合势能函数把边权调整为非负，以便使用 Dijkstra。

方法步骤：
\begin{enumerate}
  \item 建立残量网络，并为每条正向边添加反向边；
  \item 在残量网络中求从 $s$ 到 $t$ 的最短费用路；
  \item 沿该路增广可行流量；
  \item 更新残量容量与反向边费用；
  \item 直到无法增广，或达到指定流量。
\end{enumerate}
应用包括运输、任务分配、供需平衡、带费用的匹配等。实现时最容易错的是反向边费用、势能更新和负费用处理。

\section{二分图匹配}
二分图 $G=(L\cup R,E)$ 的边只连接左部点与右部点。匹配是一组两两不共享端点的边。最大匹配可以转化为最大流：从源点连到所有左部点，容量为 1；左部点到右部点按原边连接，容量为 1；右部点连到汇点，容量为 1。

Hopcroft-Karp 算法比逐条增广更快。它每一轮先 BFS 建立“最短增广路层次图”，再 DFS 同时寻找一组点不相交的最短增广路，从而把复杂度降到
\[
\bigO(E\sqrt V).
\]

\begin{summarybox}
MST 的核心是割性质：Kruskal 按边从小到大，Prim 按顶点逐步扩张。最大流的核心是残量网络，定理终点是最大流等于最小割。最小费用最大流是在流量基础上叠加费用优化。二分图匹配既可用最大流建模，也可用 Hopcroft-Karp 高效求解。
\end{summarybox}

\chapter{组合优化：TSP、图着色与启发式方法}
\section{旅行商问题 TSP}
\subsection{问题与边界}
TSP 要求从一个城市出发，经过每个城市恰好一次后回到起点，并使总路程最短。它是经典 NP-困难问题。小规模可用动态规划精确求解，大规模通常用启发式或近似方法。

Held-Karp 动态规划设 $dp[S][j]$ 表示从起点 $0$ 出发，访问集合 $S$ 且最后停在 $j$ 的最短距离：
\[
dp[S][j]=\min_{i\in S, i\ne j}\{dp[S\setminus\{j\}][i]+d(i,j)\}.
\]
最终答案为
\[
\min_{j\ne0}\{dp[V][j]+d(j,0)\}.
\]
复杂度约为 $\bigO(n^2 2^n)$，只能用于较小 $n$。

\begin{figure}[htbp]
\centering
\begin{tikzpicture}[vertex/.style={circle,draw,minimum size=7mm},>=Latex]
\node[vertex] (A) at (0,0) {1};
\node[vertex] (B) at (2,1.6) {2};
\node[vertex] (C) at (5,1.2) {3};
\node[vertex] (D) at (4,-1.3) {4};
\node[vertex] (E) at (1.3,-1.6) {5};
\draw[thick,->] (A)--(B);\draw[thick,->] (B)--(C);\draw[thick,->] (C)--(D);\draw[thick,->] (D)--(E);\draw[thick,->] (E)--(A);
\end{tikzpicture}
\caption{TSP 巡回路径示意}
\end{figure}

\section{图着色}
图着色要求给每个顶点分配颜色，使相邻顶点颜色不同；最少颜色数称为色数 $\chi(G)$。应用包括排课、寄存器分配、无线频率分配。一般图最优着色是难问题，因此常用贪心、回溯、DSATUR 等方法。

贪心着色步骤：按照某种顶点顺序，依次给当前顶点选择最小可用颜色。它快，但不保证最优。DSATUR 的思想是优先处理“饱和度”高的顶点，即邻居中已经出现颜色种类多的顶点。

\section{遗传算法 GA}
遗传算法把一个候选解编码为染色体。对 TSP，可把城市排列作为染色体。典型流程如下：
\begin{enumerate}
  \item 初始化若干随机个体；
  \item 计算适应度，例如路径越短适应度越高；
  \item 选择较优个体作为父代；
  \item 交叉生成子代；
  \item 小概率变异以增加多样性；
  \item 重复迭代并保留最好解。
\end{enumerate}
GA 的优点是模型灵活，缺点是参数敏感、收敛不稳定，不能保证全局最优。

\section{模拟退火 SA}
模拟退火从一个解出发，在邻域中随机移动。若新解更好则接受；若更差，也以概率
\[
P=\exp\left(-\frac{\Delta}{T}\right)
\]
接受，其中 $\Delta$ 是变差的幅度，$T$ 是温度。温度高时敢于跳出局部最优，温度低时趋向稳定收敛。

用于 TSP 时，常用 2-opt 作为邻域：选两条边断开并反转中间路径。模拟退火的关键不是公式，而是温度初值、降温速度、每个温度下的迭代次数。

\section{蚁群算法 ACO}
蚁群算法模拟蚂蚁利用信息素寻找路径。对 TSP，蚂蚁从一个城市出发，按“信息素强度”和“启发式距离”选择下一个城市。常见转移概率为
\[
P_{ij}=\frac{\tau_{ij}^{\alpha}\eta_{ij}^{\beta}}{\sum_{k\in allowed}\tau_{ik}^{\alpha}\eta_{ik}^{\beta}},
\]
其中 $\tau_{ij}$ 是信息素，$\eta_{ij}=1/d_{ij}$ 是启发信息。每轮后，信息素挥发并在较优路径上增加。

\begin{longtable}{p{2.5cm}p{3.2cm}p{4.4cm}p{4.2cm}}
\caption{组合优化方法对比}\\
\toprule
方法 & 结果性质 & 优点 & 局限 \\
\midrule
\endfirsthead
\toprule
方法 & 结果性质 & 优点 & 局限 \\
\midrule
\endhead
Held-Karp & 精确最优 & 理论清晰，小规模可靠 & 指数复杂度，规模稍大不可用 \\
贪心着色 & 近似/启发式 & 简单快速 & 强依赖顶点顺序，不保最优 \\
遗传算法 & 启发式 & 适合编码清晰的大规模问题 & 参数敏感，可能早熟 \\
模拟退火 & 启发式 & 能跳出局部最优，容易实现 & 降温策略影响很大 \\
蚁群算法 & 启发式 & 适合路径类组合优化 & 信息素参数多，易陷入早熟 \\
\bottomrule
\end{longtable}

\begin{summarybox}
TSP 与图着色要先判断“是否必须最优”。若规模小，可用动态规划或回溯；若规模大，通常选 2-opt、GA、SA、ACO 等启发式。启发式的正确说法是“给出较好可行解”，不要把它说成保证最优。
\end{summarybox}

\chapter{树结构：B树、B+树、红黑树与赫夫曼树}
\section{B树}
\subsection{导言}
B树是面向外存和数据库索引的多路平衡搜索树。二叉搜索树每个节点最多两个孩子，若数据在磁盘中，树高过大就会造成大量 I/O。B树让一个节点保存多个关键字与多个孩子，使树的高度显著降低。

一棵 $m$ 阶 B树通常满足：
\begin{itemize}
  \item 每个节点最多有 $m$ 个孩子；
  \item 除根外的内部节点至少有 $\lceil m/2\rceil$ 个孩子；
  \item 含 $k$ 个孩子的节点有 $k-1$ 个关键字；
  \item 所有叶子在同一层。
\end{itemize}

\subsection{查找、插入、分裂}
查找与普通搜索树类似：在节点内部二分查找关键字区间，再进入对应孩子。插入时先插入叶节点；若节点关键字过多，就把中间关键字上移，左右两部分分裂成两个节点。分裂可能向上传播，甚至产生新根。

\begin{figure}[htbp]
\centering
\begin{tikzpicture}[node distance=10mm,every node/.style={draw,rectangle,minimum height=7mm,align=center},level distance=12mm]
\node (r) at (0,0) {$17\mid35$};
\node (a) at (-3,-1.5) {$5\mid9\mid12$};
\node (b) at (0,-1.5) {$20\mid28$};
\node (c) at (3,-1.5) {$40\mid48\mid60$};
\draw (r)--(a);\draw (r)--(b);\draw (r)--(c);
\node[draw=none] at (0,0.85) {B树：内部节点保存关键字，也可保存记录指针};
\end{tikzpicture}
\caption{B树页结构示意}
\end{figure}

\section{B+树}
B+树是数据库索引中更常见的结构。与 B树相比，B+树通常让内部节点只保存索引关键字，不保存完整记录；所有实际记录或记录指针集中在叶节点，叶节点之间还按顺序相连。因此，B+树非常适合范围查询：找到范围起点后，沿叶子链表向右扫描即可。

\begin{figure}[htbp]
\centering
\begin{tikzpicture}[every node/.style={draw,rectangle,minimum height=7mm,align=center},>=Latex]
\node (r) at (0,0) {$17\mid35$};
\node (a) at (-3,-1.5) {$5\mid9\mid12$};
\node (b) at (0,-1.5) {$17\mid20\mid28$};
\node (c) at (3,-1.5) {$35\mid40\mid48$};
\draw (r)--(a);\draw (r)--(b);\draw (r)--(c);
\draw[->,thick] (a.east)--(b.west);\draw[->,thick] (b.east)--(c.west);
\node[draw=none] at (0,0.85) {B+树：内部节点导航，叶节点链表支持范围查询};
\end{tikzpicture}
\caption{B+树索引与叶子链表}
\end{figure}

\section{红黑树}
红黑树是自平衡二叉搜索树。它给每个节点增加红/黑颜色，并通过旋转与变色维持近似平衡。常用性质如下：
\begin{enumerate}
  \item 每个节点为红色或黑色；
  \item 根节点为黑色；
  \item 所有空叶子 NIL 视为黑色；
  \item 红色节点不能有红色孩子；
  \item 从任一节点到其所有后代 NIL 叶子的简单路径，包含相同数量的黑节点。
\end{enumerate}
这些性质保证树高为 $\bigO(\log n)$，因此查找、插入、删除均为 $\bigO(\log n)$。

\subsection{旋转的直觉}
旋转不破坏二叉搜索树的中序顺序，只改变局部高度。左旋可理解为“让右孩子上升，当前节点下降到其左侧”；右旋相反。

\begin{figure}[htbp]
\centering
\begin{tikzpicture}[node distance=12mm,vertex/.style={circle,draw,minimum size=7mm},>=Latex]
\node[vertex] (x) at (0,0) {$x$};
\node[vertex] (a) at (-1,-1.1) {$A$};
\node[vertex] (y) at (1,-1.1) {$y$};
\node[vertex] (b) at (0.4,-2.2) {$B$};
\node[vertex] (c) at (1.6,-2.2) {$C$};
\draw (x)--(a);\draw (x)--(y);\draw (y)--(b);\draw (y)--(c);
\node[draw=none] at (2.7,-1) {$\xrightarrow{\text{左旋 }x}$};
\node[vertex] (y2) at (4.5,0) {$y$};
\node[vertex] (x2) at (3.5,-1.1) {$x$};
\node[vertex] (c2) at (5.5,-1.1) {$C$};
\node[vertex] (a2) at (2.9,-2.2) {$A$};
\node[vertex] (b2) at (4.1,-2.2) {$B$};
\draw (y2)--(x2);\draw (y2)--(c2);\draw (x2)--(a2);\draw (x2)--(b2);
\end{tikzpicture}
\caption{红黑树左旋保持中序顺序：A < x < B < y < C}
\end{figure}

插入修复的基本思路：新节点先染红；若父亲是黑色则结束；若父亲与叔叔都是红色，则父叔变黑、祖父变红并继续向上；若叔叔是黑色，则通过旋转把“连续红色”修复掉。

\section{赫夫曼树}
赫夫曼树用于构造最优前缀编码。前缀编码要求任一字符编码都不是另一字符编码的前缀，这样解码时不会产生歧义。若字符频率不同，高频字符应使用短编码，低频字符使用长编码。

赫夫曼算法步骤：
\begin{enumerate}
  \item 把每个字符看成一棵权为频率的单节点树；
  \item 每次取权最小的两棵树合并，生成新树，权为二者之和；
  \item 重复直到只剩一棵树；
  \item 左边记 0，右边记 1，即可得到编码。
\end{enumerate}

\begin{figure}[htbp]
\centering
\begin{tikzpicture}[level distance=11mm,sibling distance=24mm,every node/.style={circle,draw,minimum size=7mm,align=center},edge from parent/.style={draw,-Latex}]
\node {100}
 child {node {45\\A}}
 child {node {55}
   child {node {25\\B}}
   child {node {30}
     child {node {12\\C}}
     child {node {18\\D}}
   }
 };
\end{tikzpicture}
\caption{赫夫曼树示意：权值越大离根越近}
\end{figure}

\begin{lstlisting}
import heapq
from itertools import count

def huffman(freq):
    uid = count()
    heap = [(w, next(uid), ch) for ch, w in freq.items()]
    heapq.heapify(heap)
    while len(heap) > 1:
        w1, _, t1 = heapq.heappop(heap)
        w2, _, t2 = heapq.heappop(heap)
        heapq.heappush(heap, (w1 + w2, next(uid), (t1, t2)))
    tree = heap[0][2]
    codes = {}
    def walk(t, code):
        if isinstance(t, tuple):
            walk(t[0], code + "0")
            walk(t[1], code + "1")
        else:
            codes[t] = code or "0"
    walk(tree, "")
    return codes
\end{lstlisting}

\begin{summarybox}
B树和 B+树是外存索引结构，重点是降低树高和减少磁盘 I/O；红黑树是内存中的有序字典结构，重点是用颜色和旋转保持对数高度；赫夫曼树是编码树，重点是“每次合并最小权”的贪心最优性。
\end{summarybox}

\chapter{全局比较、建模套路与复习结论}
\section{方法总表}
\begin{longtable}{p{2.7cm}p{3.2cm}p{3.3cm}p{4.7cm}}
\caption{图论、网络算法与树结构总览}\\
\toprule
方法 & 典型复杂度 & 适用条件 & 关键结论/限制 \\
\midrule
\endfirsthead
\toprule
方法 & 典型复杂度 & 适用条件 & 关键结论/限制 \\
\midrule
\endhead
DFS & $\bigO(V+E)$ & 图遍历、连通性、环检测 & 深递归可能爆栈，可改显式栈 \\
BFS & $\bigO(V+E)$ & 无权最短路、层次扩展 & 不适用于带权最短路 \\
拓扑排序 & $\bigO(V+E)$ & DAG 依赖关系 & 有环则不存在拓扑序 \\
关键路径 & $\bigO(V+E)$ & 项目网络 DAG & 本质是 DAG 最长路 \\
Dijkstra & $\bigO((V+E)\log V)$ & 非负权单源最短路 & 负权边会破坏贪心 \\
Bellman-Ford & $\bigO(VE)$ & 可含负边 & 可检测负环，速度较慢 \\
Floyd-Warshall & $\bigO(V^3)$ & 全源最短路 & 适合点数较小或稠密图 \\
Johnson & 约 $\bigO(VE\log V)$ & 稀疏图全源，可有负边 & 不能有负环，依赖重赋权 \\
A* & 依赖启发函数 & 有目标点路径搜索 & 启发函数越好，搜索越少 \\
Kruskal & $\bigO(E\log E)$ & 无向连通图 MST & 稀疏图常用，需并查集 \\
Prim & $\bigO(E\log V)$ & 无向连通图 MST & 稠密图可用矩阵版 \\
最大流 & Edmonds-Karp $\bigO(VE^2)$ & 容量网络 & 最大流值等于最小割容量 \\
最小费用流 & 依赖流量与最短路实现 & 容量加费用网络 & 注意反向边与负费用 \\
二分图匹配 & $\bigO(E\sqrt V)$ & 二分图 & 可转最大流，也可 Hopcroft-Karp \\
TSP & $\bigO(n^2 2^n)$ 精确 DP & 小规模精确 & 大规模通常启发式 \\
图着色 & 一般难问题 & 冲突分配、排课 & 贪心不保最优 \\
GA & 迭代式 & 大规模近似优化 & 参数敏感，需防早熟 \\
SA & 迭代式 & 局部搜索问题 & 降温策略关键 \\
ACO & 迭代式 & 路径组合优化 & 信息素参数敏感 \\
B树/B+树 & $\bigO(\log n)$ & 外存索引 & B+树范围查询更方便 \\
红黑树 & $\bigO(\log n)$ & 内存有序映射 & 旋转与变色维护平衡 \\
赫夫曼树 & $\bigO(n\log n)$ & 前缀编码 & 高频短码、低频长码 \\
\bottomrule
\end{longtable}

\section{建模套路}
遇到题目时，可按以下顺序判断：
\begin{enumerate}
  \item 是否有“先后依赖”？若有，先想 DAG、拓扑排序、关键路径。
  \item 是否要求“最短、最省、最小代价路径”？看边权：无权 BFS，非负 Dijkstra，有负边 Bellman-Ford，全源 Floyd 或 Johnson。
  \item 是否要求“连接所有点且总成本最小”？用 MST，边排序明显时 Kruskal，点扩张明显时 Prim。
  \item 是否出现“容量、供给、需求、最多能运多少”？用最大流；若同时出现费用，用最小费用流。
  \item 是否是“人-任务、课程-教室、男-女、左-右”两类对象配对？优先想到二分图匹配。
  \item 是否要求“访问所有点一次、排课不冲突、最少颜色”？通常是难问题，要说明精确算法规模限制与启发式方案。
  \item 是否是“数据库索引、磁盘页、范围查询”？考虑 B树/B+树；若是内存有序字典，考虑红黑树；若是压缩编码，考虑赫夫曼树。
\end{enumerate}

\section{最后结论}
这些算法不是孤立记忆的。它们可以按三条线串起来：第一条是“搜索线”，从 DFS/BFS 到拓扑排序、最短路和 A*；第二条是“优化线”，从 MST 到最大流、最小费用流、匹配、TSP 和图着色；第三条是“数据结构线”，从多路外存索引到内存平衡树，再到编码树。复习时应抓住每个方法的三个问题：它解决什么模型，成立条件是什么，为什么这个步骤正确。代码只是把这些思想落地。

\backmatter
\chapter*{参考来源与延伸阅读}
\addcontentsline{toc}{chapter}{参考来源与延伸阅读}
本文主要依据经典论文、算法词典与 LaTeX 官方资料整理。下列链接适合继续查阅原始定义、历史来源和排版工具文档。
\begin{enumerate}
\item E. W. Dijkstra, \textit{A Note on Two Problems in Connexion with Graphs}, 1959. \url{https://link.springer.com/article/10.1007/BF01386390}
\item L. R. Ford and D. R. Fulkerson, \textit{Maximal Flow Through a Network}, 1956. \url{https://www.cs.yale.edu/homes/lans/readings/routing/ford-max_flow-1956.pdf}
\item J. E. Hopcroft and R. M. Karp, \textit{An $n^{5/2}$ Algorithm for Maximum Matchings in Bipartite Graphs}, 1973. \url{https://epubs.siam.org/doi/10.1137/0202019}
\item P. E. Hart, N. J. Nilsson and B. Raphael, \textit{A Formal Basis for the Heuristic Determination of Minimum Cost Paths}, 1968. \url{https://ieeexplore.ieee.org/document/4082128/}
\item D. A. Huffman, \textit{A Method for the Construction of Minimum-Redundancy Codes}, 1952. \url{https://ieeexplore.ieee.org/document/4051119}
\item R. Bayer and E. M. McCreight, \textit{Organization and Maintenance of Large Ordered Indexes}, 1972. \url{https://link.springer.com/article/10.1007/BF00288683}
\item L. J. Guibas and R. Sedgewick, \textit{A Dichromatic Framework for Balanced Trees}, 1978. \url{https://sedgewick.io/wp-content/themes/sedgewick/papers/1978Dichromatic.pdf}
\item NIST Dictionary of Algorithms and Data Structures. \url{https://www.nist.gov/dads/}
\item CTAN ctex: LaTeX classes and packages for Chinese typesetting. \url{https://ctan.org/pkg/ctex?lang=en}
\item CTAN PGF/TikZ: Create PostScript and PDF graphics in TeX. \url{https://ctan.org/pkg/pgf}
\item CTAN booktabs: Publication quality tables in LaTeX. \url{https://ctan.org/pkg/booktabs?lang=en}
\end{enumerate}

\end{document}
